\documentclass[10pt,a4paper]{amsart}
\usepackage[margin=19mm]{geometry}
\usepackage{amsmath,amssymb,mathtools}
\usepackage[scr=rsfs,cal=euler]{mathalfa}
\usepackage{xfrac}
\usepackage[unicode,hidelinks,bookmarks=false]{hyperref}
\newcommand{\ie}{i.e.\ }
\newcommand{\cf}{cf.\ }
\newcommand{\resp}{respectively\ }
\newcommand{\Subsection}{\subsection}
\providecommand{\nobreakdash}{\nobreak\hbox{-}}
\setlength{\emergencystretch}{2em}
\multlinegap0pt
\usepackage{bm}
\usepackage{bbm}
\usepackage{dsfont}
\usepackage{xspace}
\usepackage{cancel}
\usepackage{mathtools}
\usepackage{amsthm}
\makeatletter
\@ifundefined{lemma}{\newtheorem{lemma}{Lemma}}{}
\@ifundefined{theorem}{\newtheorem{theorem}[lemma]{Theorem}}{}
\makeatother
\newcommand{\Grad}{\nabla}
\newcommand{\T}{\mathcal T}
\renewcommand{\P}{\mathcal P}
\title[A Convergent Finite Element Scheme for the Q-Tensor Model of]{A Convergent Finite Element Scheme for the Q-Tensor Model of Liquid Crystals Subjected to an Electric Field (Theorem 2)}
\author{Max Hirsch, Franziska Weber}
\date{Source: arXiv:2307.11229v3 (CC BY 4.0)}
\begin{document}
\maketitle

\noindent\emph{Source and licence.} A Convergent Finite Element Scheme for the Q-Tensor Model of Liquid Crystals Subjected to an Electric Field. Version arXiv:2307.11229v3 (CC BY 4.0), distributed under the Creative Commons Attribution 4.0 International licence, as recorded in the arXiv licence metadata for this exact version and shown on the abstract page https://arxiv.org/abs/2307.11229v3. Full rights notice: licenses/arxiv-2307.11229-CC-BY-4.0.txt.\par\smallskip

\noindent\textit{Editorial scope.} Counted unit: the Theorem printed as Theorem 2 of the article (source0.tex lines 435--437, source label \texttt{None}), delivered together with the article's own complete original proof of it (source0.tex lines 438--514, one \texttt{proof} environment of 77 lines). No mathematics is written, repaired or re-derived: statement and proof are the article's own text line for line. Required context retained uncounted: eq:fully\_discrete\_scheme\_definition (lines 375--387); eq:fixed\_point\_Q\_equation (lines 411--421). Every definition, display and label of those lines is retained unchanged. Omitted from this package and not redistributed: 1--374, 388--410, 422--434, 515--1583. Nothing inside a retained excerpt is omitted or altered: every retained body is delivered exactly as the article's lines are written, so no omission receipt of any kind accompanies this package. Citations: the retained excerpts carry no citation command at all, so no bibliography entry of the article is transcribed and no reference is redistributed. Reference adaptation: none was needed and none was made; every reference inside the delivered unit resolves inside it, so no unresolved reference or citation is produced. Printed slips kept literal: no printed word, symbol, hypothesis or number of the counted statement or of its proof was altered. Layout and class: the article's own class is \texttt{amsart} with packages including color, multicol, empheq, geometry, subfig, listings, amsmath, amsthm, amssymb, wasysym, verbatim, bbm, graphics, url; the wrapper is the lane's pinned \texttt{amsart} editorial wrapper at 10pt on A4, which restates the theorem-like environments the retained text needs and adds the article's own macro definitions that the retained text needs (\texttt{\textbackslash Grad}, \texttt{\textbackslash P}, \texttt{\textbackslash T}), with extra wrapper packages (bm, bbm, dsfont, xspace, cancel, mathtools). The delivered numbering is the unit's own. Assets: no retained span contains a figure environment, a graphics-inclusion command, a PDF-page inclusion, a table or a TikZ picture; no image file is copied into this package, the package declares no asset dependency and no image file is redistributed with it. Licence: version arXiv:2307.11229v3 of this article is distributed under the Creative Commons Attribution 4.0 International licence, as recorded in the arXiv licence metadata for this exact version and shown on the abstract page https://arxiv.org/abs/2307.11229v3. A full rights notice is delivered at licenses/arxiv-2307.11229-CC-BY-4.0.txt. Characters: every retained excerpt is pure ASCII, so the delivered engine, XeLaTeX with its default Latin Modern text font, reports no missing character.
\begin{subequations}
	\label{eq:fully_discrete_scheme_definition}
	\begin{equation}\label{eq:Qfullydiscrete}
	\begin{split}
	- \int_\Omega D_t^+ Q_h^n : \Phi_h\, dx &= L \int_\Omega \nabla Q_h^{n+1/2}:\nabla\Phi_h\, dx +  \int_\Omega \left(\frac{\partial\mathcal F_1(Q_h^{n+1})}{\partial Q} - \frac{\partial\mathcal F_2(Q_h^n)}{\partial Q}\right):\Phi_h\, dx\\ 
	&- \frac{\varepsilon_2}{2} \int_\Omega\left((\tilde\P_h^n\odot \nabla u_h^n(\nabla u_h^{n+1})^\top)_S - \frac1d\operatorname{tr}(\tilde\P_h^n\odot \nabla u_h^n(\nabla u_h^{n+1})^\top)I\right):\Phi_h\, dx\\ 
	&- \varepsilon_3 \int_\Omega \left(\nabla u_h^{n+1/2}\cdot\operatorname{div}((\Phi_h)_S) - \frac1d\nabla u_h^{n+1/2}\cdot\nabla\operatorname{tr}\Phi_h\right)\, dx,
	\end{split}
	\end{equation}
	\begin{equation}\label{eq:ufullydiscrete}
	\int_\Omega (\varepsilon_1\nabla u_h^{n} + \varepsilon_2 \T_R(Q_h^n)\nabla u_h^{n} + \varepsilon_3\operatorname{div}(Q_h^{n}))\cdot\nabla\psi_h\, dx = 0
	\end{equation}
\end{subequations}

    \begin{equation}
    \label{eq:fixed_point_Q_equation}
    \begin{split}
		&- \int_\Omega \frac{\widehat{Q}_h^{n+1} + \tilde q_h - Q_h^n}{\Delta t} : \Phi_h\, dx\\ 
		&= L \int_\Omega \frac{\nabla\widehat{Q}_h^{n+1}+\nabla \tilde q_h+\nabla Q_h^n}{2}:\nabla\Phi_h\, dx +  \int_\Omega \left(\frac{\partial\mathcal F_1(Q_h^{n+1})}{\partial Q} - \frac{\partial\mathcal F_2(Q_h^n)}{\partial Q}\right):\Phi_h\, dx\\ 
		&- \frac{\varepsilon_2}{2} \int_\Omega\left((\tilde\P(Q_h^{n+1},Q_h^n)\odot \nabla u_h^n(\nabla \widehat{u}_h^{n+1})^\top)_S - \frac1d\operatorname{tr}(\tilde\P(Q_h^{n+1},Q_h^n)\odot \nabla u_h^n(\nabla \widehat{u}_h^{n+1})^\top)I\right):\Phi_h\, dx\\ 
		&- \frac{\varepsilon_2}{2} \int_\Omega\left((\tilde\P(Q_h^{n+1},Q_h^n)\odot \nabla u_h^n(\nabla \tilde g_h^{n+1})^\top)_S - \frac1d\operatorname{tr}(\tilde\P(Q_h^{n+1},Q_h^n)\odot \nabla u_h^n(\nabla \tilde g_h^{n+1})^\top)I\right):\Phi_h\, dx\\
		&- \varepsilon_3 \int_\Omega \left(\frac{\nabla \widehat{u}_h^{n+1} + \nabla u_h^n}{2}\cdot\operatorname{div}((\Phi_h)_S) - \frac1d\frac{\nabla \widehat{u}_h^{n+1}+\nabla u^n}{2}\cdot\nabla\operatorname{tr}\Phi_h\right)\, dx\\
		&- \frac{\varepsilon_3}{2} \int_\Omega \left(\nabla \tilde g_h^{n+1}\cdot\operatorname{div}((\Phi_h)_S) - \frac1d\nabla \tilde g_h^{n+1}\cdot\nabla\operatorname{tr}\Phi_h\right),
\end{split}
\end{equation}

\begin{theorem}[Trace-Free and Symmetry Preservation Generalized]
If $Q_h^n$ is trace-free and symmetric and $\lambda\in (0,1]$, we have that every $\widehat{Q}_h^{n+1}$ which satisfies $\lambda^{-1}\widehat{x} = \mathcal{L}\widehat{x}$ is trace-free and symmetric. In particular, for $\lambda = 1$, this shows that solutions $Q_h^{n+1}$ to (\ref{eq:fully_discrete_scheme_definition}) are trace-free and symmetric.
\end{theorem}

\begin{proof}
	Let $\lambda \in (0,1]$ and let $\widehat{Q}_h^{n+1}$ satisfy $\lambda^{-1}\widehat{x} = \mathcal{L}\widehat{x}$. We begin by showing that $\widehat{Q}_h^{n+1}$ is symmetric. Take $\Phi_h = \widehat{Q}_h^{n+1} - (\widehat{Q}_h^{n+1})^\top := V_h^{n+1}$ in \eqref{eq:fixed_point_Q_equation}. Note that $V_h^{n+1}$ is skew-symmetric. Writing $\overline Q_h^{n+1} = \widehat{Q}_h^{n+1} + \tilde q_h$, it follows that
	\[
	\int_\Omega\left((\tilde\P(\overline Q_h^{n+1},Q_h^n)\odot \nabla u_h^n(\nabla \widehat{u}_h^{n+1})^\top)_S - \frac{1}{d}\operatorname{tr}(\tilde\P(\overline Q_h^{n+1},Q_h^n)\odot \nabla u_h^n(\nabla \widehat{u}_h^{n+1})^\top)I\right):V_h^{n+1}\, dx = 0
	\]
	since $(\tilde\P(\overline Q_h^{n+1},Q_h^n)\odot \nabla u_h^n(\nabla \widehat{u}_h^{n+1})^\top)_S - \frac{1}{d}\operatorname{tr}(\tilde\P(\overline Q_h^{n+1},Q_h^n)\odot \nabla u_h^n(\nabla \widehat{u}_h^{n+1})^\top)I$ is symmetric and $V_h^{n+1}$ is skew-symmetric. For the same reason, 
	\begin{equation*}
	\int_\Omega\left((\tilde\P(\overline Q_h^{n+1},Q_h^n)\odot \nabla u_h^n(\nabla \tilde g_h^{n+1})^\top)_S - \frac{1}{d}\operatorname{tr}(\tilde\P(\overline Q_h^{n+1},Q_h^n)\odot \nabla u_h^n(\nabla \tilde g_h^{n+1})^\top)I\right):V_h^{n+1}\, dx = 0.
	\end{equation*}
	Furthermore, we have
	\begin{equation*}
	\int_\Omega \left(\frac{\nabla \widehat{u}_h^{n+1}+\nabla u_h^n}{2}\cdot\operatorname{div}((V_h^{n+1})_S) - \frac{1}{d}\frac{\nabla \widehat{u}_h^{n+1}+\nabla u_h^n}{2}\cdot\nabla\operatorname{tr}V_h^{n+1}\right)\, dx = 0
	\end{equation*}
	and
	\begin{equation*}
	\int_\Omega \left(\nabla\tilde g_h^{n+1}\cdot \operatorname{div}((V_h^{n+1})_S) - \frac1d\nabla\tilde g_h^{n+1}\cdot\nabla\operatorname{tr}V_h^{n+1}\right)\, dx = 0
	\end{equation*}
	since $(V_h^{n+1})_S = 0$ and $\operatorname{tr}V_h^{n+1} = 0$. We also have
	\[
	\int_\Omega \tilde q_h : V_h^{n+1}\, dx = \int_\Omega \nabla \tilde q_h : \nabla V_h^{n+1}\, dx = \int_\Omega Q_h^n:V_h^{n+1}\, dx = \int_\Omega \nabla Q_h^n : \nabla V_h^{n+1}\, dx = 0
	\]
	because $\tilde q_h$ and $Q_h^n$ are symmetric and $V_h^{n+1}$ is skew-symmetric. Using these facts again gives
	\[
	\int_\Omega \frac{\partial \mathcal{F}_2(Q_h^n)}{\partial Q}:V_h^{n+1}\, dx = \int_\Omega \left((\beta_1 - a)Q_h^n + (\beta_2 - c)\operatorname{tr}((Q_h^n)^2)Q_h^n\right):V_h^{n+1}\, dx = 0.
	\]
	Thus we have
	\begin{equation}
	\frac{1}{\lambda\Delta t}\int_\Omega \widehat{Q}_h^{n+1}:V_h^{n+1}\, dx + \frac{L}{2\lambda}\int_\Omega \nabla \widehat{Q}_h^{n+1}: \nabla V_h^{n+1}\, dx + \int_\Omega \frac{\partial \mathcal{F}_1(\overline{Q}_h^{n+1})}{\partial Q}: V_h^{n+1}\, dx = 0.
	\end{equation}
	Now because $V_h^{n+1}$ is skew-symmetric and $\left(\frac{\partial \mathcal{F}_1(\overline{Q}_h^{n+1})}{\partial Q}\right)^\top = \frac{\partial \mathcal{F}_1((\overline{Q}_h^{n+1})^\top)}{\partial Q}$, we also have
	\begin{equation}
	-\frac{1}{\lambda\Delta t}\int_\Omega (\widehat{Q}_h^{n+1})^\top:V_h^{n+1}\, dx - \frac{L}{2\lambda}\int_\Omega \nabla (\widehat{Q}_h^{n+1})^\top: \nabla V_h^{n+1}\, dx - \int_\Omega \frac{\partial \mathcal{F}_1((\overline{Q}_h^{n+1})^\top)}{\partial Q}: V_h^{n+1}\, dx = 0.
	\end{equation}
	Adding the previous two equations together gives
	\begin{equation}
	\frac{1}{\lambda\Delta t}\|V_h^{n+1}\|^2 + \frac{L}{2\lambda}\|\nabla V_h^{n+1}\|^2 + \left\langle \frac{\partial \mathcal{F}_1(\overline{Q}_h^{n+1})}{\partial Q} - \frac{\partial \mathcal{F}_1((\overline{Q}_h^{n+1})^\top)}{\partial Q}, \overline{Q}_h^{n+1} - (\overline{Q}_h^{n+1})^\top\right\rangle = 0,
	\end{equation}
 where we have used $V_h^{n+1} = \widehat{Q}_h^{n+1} - (\widehat{Q}_h^{n+1})^\top = \overline{Q}_h^{n+1} - (\overline{Q}_h^{n+1})^\top$ since $\tilde q_h^{n+1}$ is symmetric.
	And since $\mathcal{F}_1$ is convex, it follows that
	\begin{align*}
	&\left\langle \frac{\partial \mathcal{F}_1(\overline{Q}_h^{n+1})}{\partial Q} - \frac{\partial \mathcal{F}_1((\overline{Q}_h^{n+1})^\top)}{\partial Q}, \overline{Q}_h^{n+1} - (\overline{Q}_h^{n+1})^\top\right\rangle \\
	&= \left\langle \frac{\partial \mathcal{F}_1(\overline{Q}_h^{n+1})}{\partial Q} , \overline{Q}_h^{n+1} - (\overline{Q}_h^{n+1})^\top\right\rangle - \left\langle \frac{\partial \mathcal{F}_1((\overline{Q}_h^{n+1})^\top)}{\partial Q}, \overline{Q}_h^{n+1} - (\overline{Q}_h^{n+1})^\top\right\rangle\\
	&\ge -(\mathcal{F}_1((\overline{Q}_h^{n+1})^\top) - \mathcal{F}_1(\overline{Q}_h^{n+1})) - (\mathcal{F}_1(\overline{Q}_h^{n+1}) - \mathcal{F}_1((\overline{Q}_h^{n+1})^\top))\\
	&= 0.
	\end{align*}
	Thus
	\[
	\frac{1}{\lambda\Delta t} \|V_h^{n+1}\|^2 + \frac{L}{2\lambda}\|\nabla V_h^{n+1}\|^2 \le 0,
	\]
	which implies that $V_h^{n+1} \equiv 0$ so that $\widehat{Q}_h^{n+1}$ is symmetric.
	
	We now show that $\widehat{Q}_h^{n+1}$ is trace-free. Taking $\Phi_h = \operatorname{tr}(\widehat{Q}_h^{n+1})I$ as a test function in~\eqref{eq:fixed_point_Q_equation}, we have
	\begin{align*}
	&-\frac{1}{\lambda\Delta t}\int_\Omega \lvert \operatorname{tr}\widehat{Q}_h^{n+1}\rvert^2\, dx\\ 
	&\hspace{5ex}= L\int_\Omega \frac{\lambda^{-1}\nabla \widehat{Q}_h^{n+1}+\nabla\tilde q_h+\nabla Q_h^n}{2} : \nabla(\operatorname{tr}(\widehat{Q}_h^{n+1})I)\, dx + \int_\Omega \left(\frac{\partial\mathcal{F}_1(\overline{Q}_h^{n+1})}{\partial Q} - \frac{\partial\mathcal{F}_2(Q_h^n)}{\partial Q}\right) : \operatorname{tr}(\widehat{Q}_h^{n+1})I\, dx\\
	&\hspace{5ex}= \frac{L}{2\lambda}\int_\Omega \lvert \nabla \operatorname{tr}(\widehat{Q}_h^{n+1})\rvert^2\, dx + \int_\Omega \left(\frac{\partial\mathcal{F}_1(\overline{Q}_h^{n+1})}{\partial Q} - \frac{\partial\mathcal{F}_2(Q_h^n)}{\partial Q}\right) : \operatorname{tr}(\widehat{Q}_h^{n+1})I\, dx,
	\end{align*}
	where for the last equality we used $\nabla A : \nabla (\operatorname{tr}(B)I) = \nabla \operatorname{tr}(A) \cdot \nabla \operatorname{tr}(B)$ and the fact that $\tilde q_h$ and $Q_h^n$ are trace-free. The remaining terms involving $\Grad u^n_h$ and $\Grad u^{n+1}_h$ cancel out due to the structure of the test function. Then computing the second integral above, we have
	\[
	\int_\Omega \frac{\partial\mathcal{F}_2(Q_h^n)}{\partial Q} : \operatorname{tr}(\widehat{Q}_h^{n+1})I\, dx = \int_\Omega \left((\beta_1 - a)Q_h^n + (\beta_2 - c)\operatorname{tr}((Q_h^n)^2)Q_h^n\right) : \operatorname{tr}(\widehat{Q}_h^{n+1})I\, dx = 0
	\]
	since $\operatorname{tr}(Q_h^n) = 0$. Furthermore,
	\begin{align*}
	&\int_\Omega \frac{\partial\mathcal{F}_1(\overline{Q}_h^{n+1})}{\partial Q} : \operatorname{tr}(\widehat{Q}_h^{n+1})I\, dx \\
	&\hspace{5ex}= \int_\Omega \left(\beta_1 \overline{Q}_h^{n+1} - b((\overline{Q}_h^{n+1})^2 - \frac{1}{d}\operatorname{tr}((\overline{Q}_h^{n+1})^2)I) + \beta_2\operatorname{tr}((\overline{Q}_h^{n+1})^2)\overline{Q}_h^{n+1}\right) : \operatorname{tr}(\widehat{Q}_h^{n+1})I\, dx\\
	&\hspace{5ex}= \int_\Omega \left(\beta_1\lvert \operatorname{tr}(\overline{Q}_h^{n+1})\rvert^2 + \beta_2\operatorname{tr}((\overline{Q}_h^{n+1})^2)\lvert\operatorname{tr}(\overline{Q}_h^{n+1})\rvert^2\right)\, dx,
	\end{align*}
	where we have used $\operatorname{tr}(\widehat{Q}_h^{n+1}) = \operatorname{tr}(\overline{Q}_h^{n+1})$ since $\operatorname{tr}(\tilde q_h) = 0$. Using this fact again, it follows that
	\[
	\int_\Omega \left(\frac{1}{\lambda\Delta t} + \beta_1 + \beta_2\operatorname{tr}((\overline{Q}_h^{n+1})^2)\right)\lvert\operatorname{tr}(\overline{Q}^{n+1})\rvert^2\, dx + \frac{L}{2\lambda}\int_\Omega \lvert \nabla \operatorname{tr}(\overline{Q}_h^{n+1})\rvert^2\, dx = 0.
	\]
	Now since $\overline{Q}_h^{n+1}$ is symmetric, we have that $\operatorname{tr}((\overline{Q}_h^{n+1})^2) = \lvert \overline{Q}_h^{n+1}\rvert_F^2 \ge 0$. We also have $\lambda, \beta_1, \beta_2 > 0$. Thus,
	\[
	\frac{1}{\Delta t}\int_\Omega \lvert \operatorname{tr}(\overline{Q}_h^{n+1})\rvert^2\, dx = \frac{1}{\Delta t}\int_\Omega \lvert \operatorname{tr}(\widehat{Q}_h^{n+1})\rvert^2\, dx \le 0,
	\]
	which shows that $\operatorname{tr}(\widehat{Q}_h^{n+1}) \equiv 0$, as desired.
\end{proof}

\end{document}
